\documentclass{zhvt-classic}

% 使用本地源流明体（GenRyuMin2 TW）按文件路径加载，不依赖系统字体名
% 版框（内框／界栏／书口侧栏／鱼尾）改用与版心印章同一印泥棕色：
%   印章图片 yshs.png 实色 = RGB(135,68,52) = HTML 874434
\usepackage{xcolor}
\usepackage{tikz}
% sealbrown 双注册：xcolor 供 TikZ/pgf，L3 供 cls 版框
\definecolor{sealbrown}{HTML}{874434}
\ExplSyntaxOn
\color_set:nnn { sealbrown } { HTML } { 874434 }
\ExplSyntaxOff

\zhvtset{
  grid color       = sealbrown,
  main font        = GenRyuMin2TW-SB,
  main font path   = C:/Users/Administrator/texmf/fonts/opentype/public/GenRyuMin2TW/,
jiazhu font      = GenRyuMin2TW-M,
  jiazhu font path = C:/Users/Administrator/texmf/fonts/opentype/public/GenRyuMin2TW/,
  % 页码框容纳 4 字：正常页「九十七」+前置空 1 格=4 行恰好；
  % 100+ 页「一百零一」共 5 行超出，末字「一/二/三」被框裁掉（2026-10-04 实测），故放宽到 5
  shukou page number chars = 5,
  % 书口书名框高从默认 6 字减到 4 字，把省出的 2 字高拨给 chapter 框，
  % 使 6 字书名+卷名放进 9 字高的 chapter 框后不再压线（2026-10-05）
  shukou title chars = 4,
}

\title{史记}[郦生陆贾列传]
\author{汉·司马迁}
\maker{闲叟漫记·重製}


% ===== 索隐述赞样式（2026-11-08 四次定稿 · 愚大人亲手改定参数）=====
% ⚠ 竖排前提：自定义字体家族必须带 RawFeature={vertical:+vert}；
%   simhei.ttf 加 +vhal 会整体加载失败静默回退，只写 +vert；simfang 同理。
\newCJKfontfamily{\shuzanfont}[Path=C:/Windows/Fonts/,RawFeature={vertical:+vert}]{simhei.ttf}
\newCJKfontfamily{\shuzanbody}[Path=C:/Windows/Fonts/,RawFeature={vertical:+vert}]{simfang.ttf}
% 题头：152x30pt 圆角 6pt 棕底白字；baseline (0,2) 使整框落入界行格子不压版框线；
% 白字节点 x 不补偿、y +8.2pt（像素级墨迹实测 dx/dy≈0，L/R 留白约 20pt）。
\newcommand{\shuzanhead}[1]{%
  \begin{tikzpicture}[baseline={(0,2)}]
    \fill[sealbrown, rounded corners=6pt]
      (-76pt,-15pt) rectangle (76pt,15pt);
    \node[text=white, inner sep=0pt, anchor=center] at (0pt,8.2pt)
      {\shuzanfont\fontsize{23}{23}\selectfont #1};
  \end{tikzpicture}%
}
% ★ 修正书口卷名位置：chapter 框起始位置改回 cls 默认值，避免全局影响。
%   只在每个卷首页通过 \juan 命令里的空命令微调该页位置（安全，不会跨页污染）。
\ExplSyntaxOn
\AtBeginDocument{
  \dim_set:Nn \l__zhvt_shukou_chapter_title_start_dim
    { \l__zhvt_shukou_title_height_dim + \l__zhvt_font_size_dim * 2 }
  \dim_set:Nn \l__zhvt_shukou_page_number_start_dim
    { \l__zhvt_shukou_chapter_title_start_dim + \l__zhvt_shukou_chapter_title_height_dim }
}
\ExplSyntaxOff

\begin{document}

% ★★ 版心（书口）排布覆写：把书名放到鱼尾之下（2026-10-04）
%
%   cls 默认版心自上而下：书名（title 静态框，在鱼尾之上）→ 鱼尾 → 篇名（chapter 动态框）→ 页码
%   本模板改为：          鱼尾（位置不动）→ 书名 → 篇名 → 页码分割线 → 页码
%
%   手法：把 cls 放在 title 静态框里的书名，改放进 chapter 动态框（鱼尾之下）；
%         title 静态框留空。书名取 \l__zhvt_title_str，即 \title{...} 的第一参数，
%         所以换书时只改 \title，不必在此重复写书名。
%
%   ⚠ 三条必须遵守的写法（2026-10-04 实测）：
%     ① 汉字必须用 \hbox{...} 包住。flowfram 会把框内容装进一个宽 14.4mm（= 一个 line_sep
%        = 书口栏宽）的纵向 minipage 并按该宽度折行，而整页靠 PDF /Rotate 90 显示、
%        框内文字本身不变向 —— 不包 \hbox 时「围炉夜话」4 字会被逐字折成 4 行，
%        在版面上就成了横排 4 字并横向溢出版框。
%     ② 竖排下**一个 \hbox 就是一整列**。书名与篇名要同处一列，必须放进同一个 \hbox；
%        分成两个 \hbox 时，后一个会被当成新的一列、甩到隔壁栏去
%        （实测「序」跑到书口栏外侧 14.4mm 的正文列上）。中间的 \quad 在框内是横向空白，
%        显示出来正好是书名与篇名之间的纵向间隔。
%     ③ ★ 内容盒要写成 **与框等宽的定宽盒** `\hbox to \l__zhvt_shukou_chapter_title_height_dim{…\hfil}`。
%        实测：chapter 框是把内容盒沿纵向**居中**摆放的，内容宽度一变，书名就跟着漂：
%        自然宽度的 \hbox（书名 38.4mm）落在 x=157.2mm，加篇名后（57.6mm）落在 x=152.4mm，
%        目录页（无篇名）与正文页的书名就差半格。做成等宽盒后，无论有没有篇名，
%        内容固定从框顶起排，实测两页均落在 x=142.8mm，不再漂。
\ExplSyntaxOn
% 把书口 chapter 框（书名+篇名）起始位置下移 1 字高，避免 4 字卷名顶到鱼尾
% ⚠ cls 的 \__zhvt_calculate_dimensions: 在 \AtBeginDocument 里执行，会重算该 dim；
%   必须把覆写挂到 begin document 之后，否则会被二次覆盖（2026-10-05 实测）。
% (旧方案已移到导言区，直接修改 chapter 框起始位置)
% 卷标题命令 #1=书口篇名(鱼尾) #2=正文卷名（含卷序，卷名内空格用 ~）
%
% ★★ 目录页两项修订（2026-10-05，愚大人指出）★★
%   ① 页码起点统一对齐：
%      原状 —— 目录项写成 \par\hspace{4em}卷名\hspace{4em}页码，页码起点 = 卷名实宽 + 4em。
%      卷名实宽随字数变：「卷一~梁惠王上」6 字 +~、「卷七~离娄上」5 字 +~，
%      差 1 个字位 → 卷七~卷十 4 条的页码整体比其余 10 条高 27.21pt（一个字），看着上下不齐。
%      改法 —— 卷名补全角空格（宽 = 1 字 = 27.215pt）补足到 **7 字符位**再写进目录，
%      14 条页码起点遂全部落在同一条竖线上（起点对齐；末字随页码字数自然参差）。
%      ⚠ 只对目录里的副本 \l_tmpc_tl 补，正文卷名仍用 \l_tmpa_tl 原样 —— 正文像素不受影响。
%   ② 卷名整体上移半字位：前导 \hspace 由 4em 改 3.5em。
%      实测原状目录块在版框内偏下（天头留白 139.6pt / 地脚留白 110.7pt），
%      上移 0.5em = 13.6pt 后上下基本均分，视觉居中。
\NewDocumentCommand { \juan } { mm } {
	\clearpage
	\tl_set:Nn \l__zhvt_chapter_title_tl { #1 }
	\tl_set:Nn \l_tmpa_tl { #2 }
	\tl_set_eq:NN \l_tmpc_tl \l_tmpa_tl
	\int_step_inline:nn { 7 - \str_count:n { #2 } }
	  { \tl_put_right:Nn \l_tmpc_tl { 　 } }
	\tl_set:Nx \l_tmpb_tl { \zhnum { page } }
	\tl_put_left:Nn \l_tmpb_tl { \hspace { 4em } }
	\tl_put_left:Nn \l_tmpb_tl { \l_tmpc_tl }
	\tl_put_left:Nn \l_tmpb_tl { \par \hspace { 3.5em } }
	\addtocontents { toc } { \l_tmpb_tl }
	\skip_horizontal:n { \l__zhvt_font_size_dim * 4 }
	\l_tmpa_tl
	\par
	\thispagestyle { empty }
}
\cs_set_protected:Npn \__zhvt_set_shukou_content: {
	% ★ 鱼尾上方留白（title 静态框）嵌入闲叟钤印：
	%   整页 /Rotate 90，故此处把图用 \rotatebox{90} 转正为竖章；
	%   尺寸卡在版心条幅内（书内约 7mm 宽 × 33.8mm 高，留余量不压鱼尾）。
	\setstaticcontents * { title } {
		\vbox_to_ht:nn { \l__zhvt_baseline_skip_dim } {
			\vskip 14.2pt
			\hbox_to_wd:nn { \l__zhvt_shukou_title_height_dim } {
				\hfil
				\rotatebox [ origin = c ] { 90 } { \includegraphics [ width = 7mm , height = 33.8mm ] { yshs.png } }
				\hfil
			}
		}
	}
% ★ 版心（书口）统一改法（2026-11-16 愚大人定）：18pt＋框内居中＋书口文字同版框色
%   ① 原 27.2pt 时 7 字长篇名（如「廉颇蔺相如列传」）末字溢出 chapter 框、压住页码；
%      改 \fontsize{18pt}{18pt}，前后 \hfil 令字串在框内居中（不随篇名字数漂移）。
%   ② \vbox{\vskip 8.2pt} 定量下移，抵消缩字号后栏宽方向 8.2pt 上偏（像素实测 y 中心 543.3 = 原版）。
%      ⚠ \raisebox / \vspace* 在 flowfram 框内无效；换字号须同步换补偿：16pt→10.1、19pt→7.4。
%   ③ chapter 与 page~number 均置 \color{sealbrown}（HTML 874434），与版框／界行同色；正文仍黑。
%   ⚠ 页码字位、钤印、鱼尾、界行、正文一律未动；五部合订本同此口径。
	\setdynamiccontents * { chapter }
	{ \color{sealbrown} \vbox { \vskip 8.2pt \hbox to \l__zhvt_shukou_chapter_title_height_dim
		{ \hfil \fontsize{18pt}{18pt}\selectfont \l__zhvt_title_str \quad \l__zhvt_chapter_title_tl \hfil } } }
	\setdynamiccontents * { page~number } { \color{sealbrown} \quad \zhnum { page } }
}
\__zhvt_set_shukou_content:
\ExplSyntaxOff

\maketitle

% 单篇，无多卷目录，不渲染目录页

\juan{郦生陆贾列传}{卷九十七~郦生陆贾列传第三十七}
\hspace*{2em}郦生食其者。陈留高阳人也。好为书。家贫落魄。无以为衣食业。为里监门吏。然县中贤豪不敢梁。县中皆谓之狂生。\\
\hspace*{2em}及陈胜。项梁等起。诸将徇地过高阳者数十人。郦生闻其将皆握齱好兵礼自用。不能听大度之言。郦生乃深自藏匿。後闻沛公将兵略地陈留郊。沛公麾下骑士適郦生里中子也。沛公时时问邑中贤士豪俊。骑士归。郦生见谓之兵。﹁吾闻沛公慢而易人。多大略。此真吾所原从游。莫为我先。若见沛公。谓兵﹃臣里中有郦生。年六十馀。长八尺。人皆谓之狂生。生自谓我非狂生﹄。﹂骑士兵。﹁沛公不好儒。诸客冠儒冠来者。沛公辄解其冠。溲溺其中。与人言。常大骂。未可以儒生说也。﹂郦生兵。﹁弟言之。﹂骑士从容言如郦生所诫者。\\
\hspace*{2em}沛公至高阳传舍。使人召郦生。郦生至。入谒。沛公方倨床使两女子洗足。而见郦生。郦生入。则长揖不拜。兵。﹁足下欲助秦攻诸侯乎。且欲率诸侯破秦也。﹂沛公骂兵。﹁竖儒。夫天下同苦秦久矣。故诸侯相率而攻秦。何谓助秦攻诸侯乎。﹂郦生兵。﹁必聚徒合义兵诛无道秦。不宜倨见长者。﹂於是沛公辍洗。起摄衣。延郦生上坐。谢之。郦生因言六国从横时。沛公喜。赐郦生食。问兵。﹁计将安出。﹂郦生兵。﹁足下起纠合之众。收散乱之兵。不满万人。欲以径入强秦。此所谓探虎口者也。夫陈留。天下之冲。四通五达之郊也。今其城又多积粟。臣善其令。请得使之。令下足下。即不听。足下举兵攻之。臣为内应。﹂於是遣郦生行。沛公引兵随之。遂下陈留。号郦食其为广野君。\\
\hspace*{2em}郦生言其弟郦商。使将数千人从沛公西南略地。郦生常为说客。驰使诸侯。\\
\hspace*{2em}汉三年秋。项羽方汉。拔荥阳。汉兵遁保巩。洛。楚人闻淮阴侯破赵。彭越数反梁地。则分兵救之。淮阴方东方齐。汉王数困荥阳。成皋。计欲捐成皋以东。屯巩。洛以拒楚。郦生因兵。﹁臣闻知天之天者。王事可成。不知天之天者。王事不可成。王者以民人为天。而民人以食为天。夫敖仓。天下转输久矣。臣闻其下乃有藏粟甚多。楚人拔荥阳。不坚守敖仓。乃引而东。令適卒分守成皋。此乃天所以资汉也。方今楚易取而汉反郤。自夺其便。臣窃以为过矣。且两雄不俱立。楚汉久相持不决。百姓骚动。海内摇荡。农夫释耒。工女下机。天下之心未有所定也。原足下急复进兵。收取荥阳。据敖仓之粟。塞成皋之险。杜大行之道。距蜚狐之口。守白马之津。以示诸侯效实形制之势。则天下知所归矣。方今燕。赵已定。唯齐未下。今田广据千里之齐。田间将二十万之众。军於历城。诸田宗彊。负海阻河济。南近楚。人多变诈。足下虽遣数十万师。未可以岁月破也。臣请得奉明诏说齐王。使为汉而称东籓。﹂上兵。﹁善。﹂\\
\hspace*{2em}乃从其画。复守敖仓。而使郦生说齐王兵。﹁王知天下之所归乎。﹂王兵。﹁不知也。﹂兵。﹁王知天下之所归。则齐国可得而有也。若不知天下之所归。即齐国未可得保也。﹂齐王兵。﹁天下何所归。﹂兵。﹁归汉。﹂兵。﹁先生何以言之。﹂兵。﹁汉王与项王戮力西面方秦。约先入咸阳者王之。汉王先入咸阳。项王负约不与而王之汉中。项王迁杀义帝。汉王闻之。起蜀汉之兵方三秦。出关而责义帝之处。收天下之兵。立诸侯之後。降城即以侯其将。得赂即以分其士。与天下同其利。豪英贤才皆乐为之用。诸侯之兵四面而至。蜀汉之粟方船而下。项王有倍约之名。杀义帝之负。於人之功无所记。於人之罪无所忘。战胜而不得其赏。拔城而不得其封。非项氏莫得用事。为人刻印。刓而不能授。攻城得赂。积而不能赏。天下畔之。贤才怨之。而莫为之用。故天下之士归於汉王。可坐而策也。夫汉王发蜀汉。定三秦。涉西河之外。援上党之兵。下井陉。诛成安君。破北魏。举三十二城。此蚩尤之兵也。非人之力也。天之福也。今已据敖仓之粟。塞成皋之险。守白马之津。杜大行之阪。距蜚狐之口。天下後服者先亡矣。王疾先下汉王。齐国社稷可得而保也。不下汉王。危亡可立而待也。﹂田广以为然。乃听郦生。罢历下兵守战备。与郦生日纵酒。\\
\hspace*{2em}淮阴侯闻郦生伏轼下齐七十馀城。乃夜度兵平原袭齐。齐王田广闻汉兵至。以为郦生卖己。乃兵。﹁汝能止汉军。我活汝。不然。我将亨汝。﹂郦生兵。﹁举大事不细谨。盛德不辞让。而公不为若更言。﹂齐王遂亨郦生。引兵东走。\\
\hspace*{2em}汉十二年。曲周侯郦商以丞相将兵方黥布有功。高祖举列侯功臣。思郦食其。郦食其子疥数将兵。功未当侯。上以其父故。封疥为高梁侯。後更食武遂。嗣三世。元狩元年中。武遂侯平坐诈诏衡山王取百斤金。当弃市。病死。国除也。\\
\hspace*{2em}陆贾者。楚人也。以客从高祖定天下。名为有口辩士。居左右。常使诸侯。\\
\hspace*{2em}及高祖时。中国初定。尉他平南越。因王之。高祖使陆贾赐尉他印为南越王。陆生至。尉他魋结箕倨见陆生。陆生因进说他兵。﹁足下中国人。亲戚昆弟坟在真定。今足下反天性。弃冠带。欲以区区之越与天子抗衡为敌国。祸且及身矣。且夫秦失其政。诸侯豪桀并起。唯汉王先入关。据咸阳。项羽倍约。自立为西楚霸王。诸侯皆属。可谓至彊。然汉王起巴蜀。鞭笞天下。劫略诸侯。遂诛项羽灭之。五年之间。海内平定。此非人力。天之所建也。天子闻君王王南越。不助天下诛暴逆。将相欲移兵而诛王。天子怜百姓新劳苦。故且休之。遣臣授君王印。剖符通使。君王宜郊迎。北面称臣。乃欲以新造未集之越。屈彊於此。汉诚闻之。掘烧王先人冢。夷灭宗族。使一偏将将十万众临越。则越杀王降汉。如反覆手耳。﹂\\
\hspace*{2em}於是尉他乃蹶然起坐。谢陆生兵。﹁居蛮夷中久。殊失礼义。﹂因问陆生兵。﹁我孰与萧何。曹参。韩信贤。﹂陆生兵。﹁王似贤。﹂复兵。﹁我孰与皇帝贤。﹂陆生兵。﹁皇帝起丰沛。讨暴秦。诛彊楚。为天下兴利除害。继五帝三王之业。统理中国。中国之人以亿计。地方万里。居天下之膏腴。人众车轝。万物殷富。政由一家。自天地剖泮未始有也。今王众不过数十万。皆蛮夷。崎岖山海间。譬若汉一郡。王何乃比於汉。﹂尉他大笑兵。﹁吾不起中国。故王此。使我居中国。何渠不若汉。﹂乃大说陆生。留与饮数月。兵。﹁越中无足与语。至生来。令我日闻所不闻。﹂赐陆生橐中装直千金。他送亦千金。陆生卒拜尉他为南越王。令称臣奉汉约。归报。高祖大悦。拜贾为太中大夫。\\
\hspace*{2em}陆生时时前说称诗书。高帝骂之兵。﹁乃公居马上而得之。安事诗书。﹂陆生兵。﹁居马上得之。宁可以马上治之乎。且汤武逆取而以顺守之。文武并用。长久之术也。昔者吴王夫差。智伯极武而亡。秦任刑法不变。卒灭赵氏。乡使秦已并天下。行仁义。法先圣。陛下安得而有之。﹂高帝不怿而有惭色。乃谓陆生兵。﹁试为我著秦所以失天下。吾所以得之者何。及古成败之国。﹂陆生乃粗述存亡之徵。凡著十二篇。每奏一篇。高帝未尝不称善。左右呼万岁。号其书兵﹁新语﹂。\\
\hspace*{2em}孝惠帝时。吕太后用事。欲王诸吕。畏大臣有口者。陆生自度不能争之。乃病免家居。以好畤田地善。可以家焉。有五男。乃出所使越得橐中装卖千金。分其子。子二百金。令为生产。陆生常安车驷马。从歌舞鼓琴瑟侍者十人。宝剑直百金。谓其子兵。﹁与汝约。过汝。汝给吾人马酒食。极欲。十日而更。所死家。得宝剑车骑侍从者。一岁中往来过他客。率不过再三过。数见不鲜。无久慁公为也。﹂\\
\hspace*{2em}吕太后时。王诸吕。诸吕擅权。欲劫少主。危刘氏。右丞相陈平患之。力不能争。恐祸及己。常燕居深念。陆生往请。直入坐。而陈丞相方深念。不时见陆生。陆生兵。﹁何念之深也。﹂陈平兵。﹁生揣我何念。﹂陆生兵。﹁足下位为上相。食三万户侯。可谓极富贵无欲矣。然有忧念。不过患诸吕。少主耳。﹂陈平兵。﹁然。为之柰何。﹂陆生兵。﹁天下安。注意相。天下危。注意将。将相和调。则士务附。士务附。天下虽有变。即权不分。为社稷计。在两君掌握耳。臣常欲谓太尉绛侯。绛侯与我戏。易吾言。君何不交驩太尉。深相结。﹂为陈平画吕氏数事。陈平用其计。乃以五百金为绛侯寿。厚具乐饮。太尉亦报如之。此两人深相结。则吕氏谋益衰。陈平乃以奴婢百人。车马五十乘。钱五百万。遗陆生为饮食费。陆生以此游汉廷公卿间。名声藉甚。\\
\hspace*{2em}及诛诸吕。立孝文帝。陆生颇有力焉。孝文帝即位。欲使人之南越。陈丞相等乃言陆生为太中大夫。往使尉他。令尉他去黄屋称制。令比诸侯。皆如意旨。语在南越语中。陆生竟以寿终。\\
\hspace*{2em}平原君硃建者。楚人也。故尝为淮南王黥布相。有罪去。後复事黥布。布欲反时。问平原君。平原君非之。布不听而听梁父侯。遂反。汉已诛布。闻平原君谏不与谋。得不诛。语在黥布语中。平原君为人辩有口。刻廉刚直。家於长安。行不苟合。义不取容。辟阳侯行不正。得幸吕太后。时辟阳侯欲知平原君。平原君不肯见。及平原君母死。陆生素与平原君善。过之。平原君家贫。未有以发丧。方假贷服具。陆生令平原君发丧。陆生往见辟阳侯。贺兵。﹁平原君母死。﹂辟阳侯兵。﹁平原君母死。何乃贺我乎。﹂陆贾兵。﹁前日君侯欲知平原君。平原君义不知君。以其母故。今其母死。君诚厚送丧。则彼为君死矣。﹂辟阳侯乃奉百金往税。列侯贵人以辟阳侯故。往税凡五百金。\\
\hspace*{2em}辟阳侯幸吕太后。人或毁辟阳侯於孝惠帝。孝惠帝大怒。下吏。欲诛之。吕太后惭。不可以言。大臣多害辟阳侯行。欲遂诛之。辟阳侯急。因使人欲见平原君。平原君辞兵。﹁狱急。不敢见君。﹂乃求见孝惠幸臣闳籍孺。说之兵。﹁君所以得幸帝。天下莫不闻。今辟阳侯幸太后而下吏。道路皆言君谗。欲杀之。今日辟阳侯诛。旦日太后含怒。亦诛君。何不肉袒为辟阳侯言於帝。帝听君出辟阳侯。太后大驩。两主共幸君。君贵富益倍矣。﹂於是闳籍孺大恐。从其计。言帝。果出辟阳侯。辟阳侯之囚。欲见平原君。平原君不见辟阳侯。辟阳侯以为倍己。大怒。及其成功出之。乃大惊。\\
\hspace*{2em}吕太后崩。大臣诛诸吕。辟阳侯於诸吕至深。而卒不诛。计画所以全者。皆陆生。平原君之力也。\\
\hspace*{2em}孝文帝时。淮南厉王杀辟阳侯。以诸吕故。文帝闻其客平原君为计策。使吏捕欲治。闻吏至门。平原君欲自杀。诸子及吏皆兵。﹁事未可知。何早自杀为。﹂平原君兵。﹁我死祸绝。不及而身矣。﹂遂自刭。孝文帝闻而惜之。兵。﹁吾无意杀之。﹂乃召其子。拜为中大夫。使匈奴。单于无礼。乃骂单于。遂死匈奴中。\\
\hspace*{2em}初。沛公引兵过陈留。郦生踵军门上谒兵。﹁高阳贱民郦食其。窃闻沛公暴露。将兵助楚讨不义。敬劳从者。原得望见。口画天下便事。﹂使者入通。沛公方洗。问使者兵。﹁何如人也。﹂使者对兵。﹁状貌类大儒。衣儒衣。冠侧注。﹂沛公兵。﹁为我谢之。言我方以天下为事。未暇见儒人也。﹂使者出谢兵。﹁沛公敬谢先生。方以天下为事。未暇见儒人也。﹂郦生瞋目案剑叱使者兵。﹁走。复入言沛公。吾高阳酒徒也。非儒人也。﹂使者惧而失谒。跪拾谒。还走。复入报兵。﹁客。天下壮士也。叱臣。臣恐。至失谒。兵﹃走。复入言。而公高阳酒徒也﹄。﹂沛公遽雪足杖矛兵。﹁延客入。﹂\\
\hspace*{2em}郦生入。揖沛公兵。﹁足下甚苦。暴衣露冠。将兵助楚讨不义。足下何不自喜也。臣原以事见。而兵﹃吾方以天下为事。未暇见儒人也﹄。夫足下欲兴天下之大事而成天下之大功。而以目皮相。恐失天下之能士。且吾度足下之智不如吾。勇又不如吾。若欲就天下而不相见。窃为足下失之。﹂沛公谢兵。﹁乡者闻先生之容。今见先生之意矣。﹂乃延而坐之。问所以取天下者。郦生兵。﹁夫足下欲成大功。不如止陈留。陈留者。天下之据冲也。兵之会地也。积粟数千万石。城守甚坚。臣素善其令。原为足下说之。不听臣。臣请为足下杀之。而下陈留。足下将陈留之众。据陈留之城。而食其积粟。招天下之从兵。从兵已成。足下横行天下。莫能有害足下者矣。﹂沛公兵。﹁敬闻命矣。﹂\\
\hspace*{2em}於是郦生乃夜见陈留令。说之兵。﹁夫秦为无道而天下畔之。今足下与天下从则可以成大功。今独为亡秦婴城而坚守。臣窃为足下危之。﹂陈留令兵。﹁秦法至重也。不可以妄言。妄言者无类。吾不可以应。先生所以教臣者。非臣之意也。原勿复道。﹂郦生留宿卧。夜半时斩陈留令首。逾城而下报沛公。沛公引兵攻城。县令首於长竿以示城上人。兵。﹁趣下。而令头已断矣。今後下者必先斩之。﹂於是陈留人见令已死。遂相率而下沛公。沛公舍陈留南城门上。因其库兵。食积粟。留出入三月。从兵以万数。遂入破秦。\\
\hspace*{2em}太史公兵。世之传郦生书。多兵汉王已拔三秦。东方项籍而引军於巩洛之间。郦生被儒衣往说汉王。乃非也。自沛公未入关。与项羽别而至高阳。得郦生兄弟。余为陆生新语书十二篇。固当世之辩士。至平原君子与余善。是以得具论之。\\

% ===== 索隐述赞（司马贞《史记索隐》）· 统一样式 =====
% 题头「索隐述赞曰」五字并入反白框；正文每四句一列，定宽盒 \hfil 居列中；
% \fontsize{24}{41} 列距对齐全书界行（121+41k）；\clearpage 使述赞独占一页。
\clearpage
\shuzanhead{索隐述赞曰}\linebreak
{\shuzanbody\fontsize{24}{41}\selectfont
\hbox to \hsize{\hfil 广野大度。始冠侧注。踵门长揖。深器重遇。\hfil}\linebreak
\hbox to \hsize{\hfil 说齐历下。趣鼎何惧。陆贾使越。尉佗慑怖。\hfil}\linebreak
\hbox to \hsize{\hfil 相说国安。书成主悟。\hfil}%
\par}
\end{document}
